\chapter{Cơ sở lý thuyết}

Chương này chỉ trình bày khái niệm \emph{đồ án thực sự dùng}. Mỗi kỹ thuật và
công nghệ đi theo bốn câu, để người đọc khỏi lẫn trang giáo trình với phần đã
cài:

\begin{itemize}
    \item \textbf{Là gì} --- nói ngắn, bằng việc đồ án cần làm, không chép định
    nghĩa bách khoa.
    \item \textbf{Vì sao dùng} --- chỗ lệch từ vựng, chỗ bịa số Điều, hoặc chỗ
    tiếng Việt không dấu mà khái niệm đó gỡ được.
    \item \textbf{Cách dùng} --- file, tham số, nút pipeline.
    \item \textbf{Khi nào} --- lúc nào khối chạy; lúc nào bị tắt hoặc bỏ qua.
\end{itemize}

Mục so sánh ở cuối chương nêu phương án đã xem nhưng không chọn.

\section{Xử lý ngôn ngữ tự nhiên trong bài toán hỏi đáp}

Bài toán của đồ án nghe đơn giản: người dùng gõ một câu tiếng Việt, máy trả lời
và chỉ ra Điều nào đang dùng. Khác hỏi đáp trên Wikipedia tiếng Anh, kho luật
Việt Nam có khung cứng (Chương, Điều, Khoản), thuật ngữ được chuẩn hóa, và một
số Điều sai có thể kéo theo tiền phạt. Vì vậy đồ án không coi đây là chatbot
tán gẫu, mà là hỏi đáp trên kho đóng: câu trả lời phải bám đoạn văn đã tìm
được~\cite{karpukhin2020dpr, lewis2020rag}.

Hai hướng thường gặp:

\begin{itemize}
    \item \textbf{Đóng sách (close-book).} Mô hình trả lời từ những gì đã được
    nén vào trọng số lúc huấn luyện. Nhanh, nhưng không chỉ được trang luật nào.
    Khi nghị định đổi, câu trả lời cũ vẫn nghe trôi.
    \item \textbf{Mở sách (open-book).} Máy tìm vài Điều liên quan, rồi mới đọc
    và viết. Đồ án theo hướng này. Nếu không tìm ra Điều, hệ thống nói thiếu
    căn cứ, không bịa đoạn khuyên chung.
\end{itemize}

Ví dụ dễ hình dung: câu ``thử việc được bao lâu'' và câu ``sa thải nhân viên
kinh doanh'' dùng chữ đời thường. Văn bản gốc viết ``thời gian thử việc'' và
``đơn phương chấm dứt hợp đồng lao động''. Máy phải vừa khớp mặt chữ (khi người
dùng gõ đúng ``Điều~25''), vừa bắt gần nghĩa (khi chữ hoàn toàn khác).

\section{Các khái niệm nền gắn với tiếng Việt và kho luật}

\subsection{Tách từ và hai lớp token}

Tiếng Việt viết cách âm tiết, không cách từ. ``Lao động'' là một từ, hai âm
tiết. Nếu BM25 đếm ``lao'' và ``động'' riêng, độ hiếm của từ bị loãng và Điều
đúng dễ bị Điều khác lấn. Đồ án tách từ bằng \texttt{underthesea} khi dựng chỉ
mục BM25. Kênh nhúng và mô hình sinh dùng cách tách của nhà cung cấp (subword).
Hai lớp token khác nhau; đồ án chấp nhận và không ép một tokenizer cho cả hai.

\boncau
{Đơn vị chữ đưa vào BM25 (từ tiếng Việt) khác đơn vị chữ đưa vào mô hình sinh
(subword).}
{Không tách từ thì BM25 yếu trên kho luật, vì nhiều cụm hai âm tiết.}
{Tách lúc dựng cache BM25. Embedding không dùng kết quả underthesea.}
{Mỗi lần lập hoặc làm mới chỉ mục chữ. Không chạy trên trang tra cứu FTS.}

\subsection{Từ phổ biến và độ hiếm}

Văn bản luật lặp ``Điều'', ``khoản'', ``quy định''. Người dùng vẫn hay gõ
``Điều~25''. Đồ án không tự soạn danh sách từ dừng riêng cho pháp lý: BM25 đã
hạ điểm những từ xuất hiện ở quá nhiều Điều (IDF). Giữ nguyên ``Điều'' giúp câu
có số hiệu không bị mất tín hiệu.

\subsection{Người tìm, máy đọc}

Người dùng DNNVV không chọn trước Chương nào để máy đọc. Hệ thống phải tự tìm
vài Điều, rồi mới viết câu~\cite{karpukhin2020dpr}. Đồ án: kênh tìm là BM25 cộng
Chroma; kênh đọc là mô hình sinh câu, không phải mô hình khoanh một đoạn kiểu
đề thi đọc hiểu. Một câu thuế thường cần hơn một Điều.

Trang \texttt{/laws} chỉ tìm và hiện chữ, không gọi kênh đọc.

\subsection{Bịa nội dung}

Mô hình vẫn có thể viết điều không có trong ngữ cảnh, hoặc không có thật. Ở
đây dạng nguy hiểm nhất là bịa số Điều hoặc bịa mức phạt. Đồ án chặn bằng ba
lớp: bắt buộc có ngữ cảnh đã tìm; lời nhắc cấm bịa; sau khi viết xong chỉ giữ
số Điều thuộc tập vừa tìm. Diễn giải sai trên Điều đúng thì ba lớp này không
bắt --- hạn chế Chương~5.

\subsection{Cửa sổ ngữ cảnh và nhiệt độ}

Mô hình nhìn các token trong một cửa sổ có hạn; chi phí tăng nhanh khi chuỗi
dài~\cite{vaswani2017}. Đó là lý do không nhét cả kho \soDieu{} Điều vào một
lời nhắc, mà cắt còn 8 Điều sau lọc. Nhiệt độ chỉnh mức ``mạo hiểm'' khi chọn
chữ tiếp theo. Đồ án để nhiệt độ 0 khi viết câu, để cùng một câu hỏi và cùng
một tập Điều cho ra lời diễn giải ổn định hơn.

Đồ án không cài mạng Transformer: gọi API đã huấn luyện. Cổng sinh
(\texttt{LLM\_MODEL}) tách khỏi cổng nhúng (\texttt{EMBEDDINGS\_MODEL}).

\subsection{Độ chính xác, độ phủ và cách đánh giá}

Trong truy hồi thông tin, độ chính xác hỏi ``trong các Điều đưa ra, bao nhiêu
Điều đúng''; độ phủ hỏi ``trong các Điều đáng ra phải tìm, máy giữ được bao
nhiêu''. Đồ án thiết kế phễu $60 \to 8$ theo hướng độ phủ: bỏ sót hại hơn lấy
thừa. Chương~4 không công bố $F_2$, MRR hay ROUGE trên bộ nhãn nội bộ --- chưa
gắn nhãn đủ lớn. Việc đánh giá sản phẩm là chạy cụm Docker và đối chiếu Điều
gốc bằng tay.

\section{Xác thực, ủy quyền và ba lớp phần mềm}

Hội thoại và hợp đồng không được nhìn chéo giữa hai công ty. Đổi giao diện không
được đụng công thức RRF. Đồ án tách ba lớp: trình bày (React), nghiệp vụ
(FastAPI, pipeline), dữ liệu (PostgreSQL, Chroma)~\cite{thanh2024}. Xác thực
trả lời ``ai đang gọi''; ủy quyền trả lời ``được làm gì''. JWT đi kèm kiểm vai
trò và cách ly theo \texttt{user\_id} / \texttt{organization\_id}. Khi đoán định
danh hội thoại người khác, API trả 404 chứ không 403, để khỏi lộ tài nguyên
đang tồn tại. HTTPS chưa bật trong Compose nội bộ.


\section{Cấu trúc văn bản quy phạm pháp luật Việt Nam}

\subsection{Luật, nghị định, thông tư không nằm trên một trang}

Câu ``công ty mười hai lao động khai GTGT theo tháng hay theo quý'' không nằm
gọn trong một Điều của Luật Quản lý thuế. Luật nói nghĩa vụ kê khai tồn tại.
Nghị định nói ngưỡng và kỳ. Thông tư nói tờ khai. Nếu kho chỉ có luật, máy trả
lời đúng nguyên tắc và sai thao tác --- kiểu sai khó phát hiện vì nghe ``có vẻ
đúng luật''. Vì vậy đồ án đưa cả ba cấp vào kho~\cite{qlthue2019, quochoi2017,
bll2019}.

\subsection{Điều là đơn vị cắt dữ liệu}

Văn bản chia Phần -- Chương -- Mục -- Điều -- Khoản -- Điểm. Đồ án cắt theo
\textbf{Điều} vì ba lý do thực dụng:

\begin{enumerate}
    \item Người làm kế toán và nhân sự dẫn chiếu ``Điều~$n$ luật~$X$'', không
    dẫn ``đoạn 400 ký tự thứ 17''.
    \item Một Điều thường đủ nghĩa để đứng trong lời nhắc mô hình. Cắt giữa câu
    làm mất chủ ngữ của khoản, dễ khiến mô hình đọc sai ai phải làm gì.
    \item Độ dài trung bình khoảng hơn một nghìn ký tự: nhét được vài Điều vào
    ngữ cảnh mà chưa đầy cửa sổ.
\end{enumerate}

Cắt cả bộ luật thành một vector thì vô dụng: máy không biết chỗ nào trả lời câu
hỏi. Cắt theo số ký tự cố định thì hợp bài xử lý ngôn ngữ chung, nhưng phá cách
dẫn chiếu của ngành luật.

\subsection{Viện dẫn chéo}

Kho gốc có trường \texttt{extra}: danh sách chỗ một Điều trỏ sang Điều khác. Đồ
án không xây đồ thị tri thức đầy đủ. Viện dẫn được lưu và hiện thành liên kết
trên trang đọc. Người đang xem Điều~44 Luật Quản lý thuế thấy đường sang Điều
nói thời hạn nộp. Khi chạy thử, người không học luật dùng chỗ ``nhảy'' này
nhiều hơn ô chat.

\section{Truy hồi theo chữ: BM25}

Khi người dùng gõ đúng cụm có trong văn bản --- ``Điều~25'', ``thử việc'',
``hóa đơn điện tử'' --- máy cần một cách chấm điểm dựa trên chữ, không dựa trên
nghĩa. Đồ án dùng BM25 cho việc đó~\cite{robertson2009bm25}. Ý tưởng: từ nào
hiếm trong kho mà xuất hiện trong Điều thì được cộng nhiều; Điều quá dài bị trừ
nhẹ để khỏi thắng chỉ vì dài.

\boncau
{Cách xếp hạng Điều theo mức khớp mặt chữ với câu hỏi, có tính độ hiếm của từ
và độ dài văn bản~\cite{robertson2009bm25}.}
{Người dùng hay gõ số hiệu. Kênh vector dễ coi ``Điều~24'' và ``Điều~25'' là
gần nhau vì chúng đứng cạnh trong cùng chương. BM25 phân biệt được hai số đó.}
{Chạy cùng lúc với Chroma trong nút \texttt{retrieve}. $k_1 = 1{,}2$,
$b = 0{,}65$. Tách từ tiếng Việt bằng \texttt{underthesea} (backend ghim
Python~3.10), kết quả tách được ghi cache.}
{Mọi câu hỏi pháp lý đi truy hồi lai. Không chạy khi ý định là \texttt{SKIP}
hoặc khi chưa có kho.}

Công thức chuẩn được giữ nguyên. $Q$ là câu hỏi, $D$ là một Điều, $f(t,D)$ là
số lần từ $t$ xuất hiện trong Điều, $|D|$ là độ dài Điều, $\mathrm{avgdl}$ là
độ dài trung bình trong kho, $N$ là số Điều, $n(t)$ là số Điều chứa $t$:

\begin{equation}
\label{eq:bm25}
\mathrm{BM25}(Q, D) = \sum_{t \in Q} \mathrm{IDF}(t) \cdot
\frac{f(t, D) \cdot (k_1 + 1)}
     {f(t, D) + k_1 \cdot \left(1 - b + b \cdot \dfrac{|D|}{\mathrm{avgdl}}\right)}
\end{equation}

\begin{equation}
\label{eq:idf}
\mathrm{IDF}(t) = \ln\left(\frac{N - n(t) + 0.5}{n(t) + 0.5} + 1\right)
\end{equation}

Hệ số $b$ giảm nhẹ so với mặc định vì độ dài Điều trong kho dao động mạnh (điều
định nghĩa ngắn, điều xử phạt dài). BM25 chịu thua khi từ vựng lệch hoàn toàn
(``sa thải'' / ``đơn phương chấm dứt'') --- lúc đó đồ án dựa vào kênh vector.

\section{Truy hồi theo nghĩa: embedding và cosine}

Kênh thứ hai biến mỗi Điều thành một dãy số (embedding). Hai dãy gần nhau theo
góc (cosine) thì nội dung gần nghĩa, dù chữ khác nhau. Đó là lý do câu ``sa
thải'' vẫn kéo được Điều nói đơn phương chấm dứt.

\boncau
{Mỗi Điều (và mỗi câu hỏi) được đổi thành một vector; độ gần đo bằng cosine.}
{Bắt cặp gần nghĩa khi mặt chữ không trùng.}
{Mô hình \texttt{gemini-embedding-001}, cắt còn 1536 chiều. Nhúng theo Điều, lô
64 bản ghi, nghỉ dần khi gặp mã 429. Đổi số chiều thì xóa Chroma và nhúng lại.}
{Lúc nạp hoặc lập lại chỉ mục (mọi Điều) và lúc hỏi (nhúng câu hoặc đoạn HyDE).
Thiếu khóa nhúng: hỏi đáp trả 503, trang tra cứu vẫn dùng được.}

\begin{equation}
\label{eq:cosine}
\mathrm{sim}(q, d) = \frac{\mathbf{v}_q \cdot \mathbf{v}_d}
                          {\lVert \mathbf{v}_q \rVert \, \lVert \mathbf{v}_d \rVert}
\end{equation}

Cosine chuẩn hóa độ dài vector, nên Điều dài không tự thắng Điều ngắn chỉ vì
có nhiều số.

\section{Gộp hai bảng xếp hạng: Reciprocal Rank Fusion}

Điểm BM25 và điểm cosine không cùng một thước. Cộng thẳng hai cột số là tự bịa
thang đo. Reciprocal Rank Fusion chỉ nhìn thứ hạng: Điều đứng hạng 1 được cộng
nhiều, hạng 20 được cộng ít, không cần biết điểm thô là 12 hay 0,83~\cite{cormack2009rrf}.

\boncau
{Cách gộp nhiều danh sách xếp hạng bằng nghịch đảo hạng~\cite{cormack2009rrf}.}
{Hai kênh trả về hai thứ tự khác nhau; cần một danh sách chung trước khi lọc
và viết câu.}
{$k=60$, trọng số nghĩa : chữ $= 3:1$, lấy 8 Điều đưa sang bước sau. Tỉ lệ 3:1
chọn sau khi thử tay: đa số câu DNNVV là câu nghĩa, nhưng BM25 vẫn kéo đúng khi
người dùng gõ ``Điều~35''.}
{Trong nút \texttt{retrieve}, sau khi hai kênh trả danh sách. Nếu một kênh rỗng
hoàn toàn thì dùng kênh còn lại.}

\begin{equation}
\label{eq:rrf}
\mathrm{RRF}(d) = \sum_{i} \frac{w_i}{k + \mathrm{rank}_i(d)}
\end{equation}

$w_i$ là trọng số kênh, $\mathrm{rank}_i(d)$ là hạng của Điều $d$ trên kênh $i$
(hạng 1 là tốt nhất). Hằng số $k$ làm giảm khoảng cách giữa hạng gần nhau, tránh
một kênh lấn hết.

\section{RAG, HyDE, lời nhắc mô hình và rủi ro bịa nguồn}

\subsection{RAG}

Mô hình ngôn ngữ lớn viết tiếng Việt trôi, nhưng không ``mở được Điều~44'' trừ
khi đoạn Điều được đưa vào lời nhắc. RAG buộc đúng việc đó: tìm trước, viết
sau~\cite{lewis2020rag}. Khi luật đổi, người vận hành nhập văn bản mới, không
bắt buộc huấn luyện lại mạng.

\boncau
{Quy trình tìm vài Điều liên quan, rồi mới viết câu trên đúng những Điều
đó~\cite{lewis2020rag}.}
{Không lấy trí nhớ trong trọng số làm nguồn luật: dễ bịa số Điều, khó cập nhật
khi văn bản đổi.}
{Bảy nút LangGraph. Sau khi viết xong, hệ thống chỉ giữ số Điều thuộc tập vừa
tìm. Nút lọc bằng mô hình giảm Điều lạc. Xếp hạng chéo tắt vì Gemini không có
cổng rerank.}
{Mỗi lượt hỏi đáp và mỗi điều khoản hợp đồng (cùng một hàm tìm). Không chạy khi
thiếu chỉ mục vector (HTTP 503).}

Rủi ro còn lại: mô hình diễn giải sai trên Điều đúng. Hàng rào số hiệu không
bắt được lỗi đó --- hạn chế nêu lại ở Chương~5.

\subsection{HyDE}

Câu hỏi của chủ doanh nghiệp thường ngắn và giọng đời thường. Văn luật dài và
dùng chữ khác. HyDE nhờ mô hình viết một đoạn ngắn ``như thể đã trả lời bằng
giọng quy phạm'', rồi dùng đoạn đó để tìm vector~\cite{gao2023hyde}. Đoạn giả
có thể sai nội dung; đồ án không hiện đoạn này cho người dùng như đáp án. Mục
đích chỉ là kéo câu hỏi lại gần giọng kho.

\boncau
{Sinh một đoạn giọng quy phạm rồi nhúng đoạn đó để tìm vector~\cite{gao2023hyde}.}
{Câu hỏi ngắn, lệch từ vựng so với văn bản gốc.}
{Mô hình viết khoảng 5--6 câu. Viết lại nhiều biến thể đang tắt để giữ độ trễ.}
{Bật trên sản phẩm trước nút tìm. Tắt được trong \texttt{config.yaml}. Lỗi HyDE
thì tìm bằng câu gốc.}

\subsection{Lời nhắc mô hình, không tinh chỉnh trọng số}

Đồ án không có GPU để tinh chỉnh mô hình. Luật đổi thì nhập JSON rẻ hơn huấn
luyện lại. Các nút chỉ được in đầu ra hẹp: ý định chỉ \texttt{NEXT} hoặc
\texttt{SKIP}; lọc chỉ PASS hoặc DROP; câu trả lời chỉ được kết luận từ căn cứ
đang có, mặc định chủ thể là DNNVV.

\subsection{Chỉ số $F_2$}

Khi thiết kế phễu tìm (lấy 60 Điều lúc gộp, còn 8 Điều lúc viết), đồ án ưu tiên
không bỏ sót Điều quyết định. Lấy thừa còn lọc được; bỏ sót thì không bước nào
cứu. Chỉ số $F_\beta$ với $\beta=2$ diễn đạt đúng thiên hướng đó, nhưng Chương~4
\emph{không} gắn một điểm $F_2$ làm kết quả sản phẩm: chưa có bộ câu gắn nhãn
nội bộ đủ lớn.

\begin{equation}
\label{eq:fbeta}
F_\beta = (1 + \beta^2) \cdot \frac{P \cdot R}{\beta^2 P + R},\qquad \beta = 2
\end{equation}

$P$ là độ chính xác (trong các Điều đưa ra, bao nhiêu Điều đúng), $R$ là độ phủ
(trong các Điều đáng ra phải tìm, máy giữ được bao nhiêu).

\section{Tìm chữ tiếng Việt trên PostgreSQL}

Trang tra cứu là đọc văn bản, không phải hỏi đáp. Người dùng hay gõ không dấu
(``thue gtgt''). PostgreSQL không có từ điển full-text tiếng Việt sẵn.
\texttt{unaccent} bỏ dấu lúc tìm; đoạn hiện trên trang vẫn là chữ gốc còn
dấu~\cite{postgresqlfts}.

\boncau
{Tìm chữ trên cột đã lập chỉ mục, không gọi mô hình~\cite{postgresqlfts}.}
{Ô tìm trên \texttt{/laws} phải trả nhanh và đọc được khi mất khóa API.}
{Cột sinh sẵn: \texttt{unaccent} bọc hàm \texttt{immutable\_unaccent}, cấu hình
\texttt{simple}. Ba tầng: số Điều, \texttt{ts\_rank}, \texttt{pg\_trgm} (ngưỡng
0,3). Bôi đậm ở React trên chữ gốc, không dùng \texttt{ts\_headline} (hàm này
trả đoạn đã mất dấu).}
{Mỗi lần gõ ô tìm. Không cần Chroma.}

\section{Công nghệ hiện thực và lý do chọn}

Bảng~\ref{tab:cong-nghe} liệt kê ngăn xếp. Mục tiếp theo viết từng mục theo đúng
bốn câu đã nêu.

\begin{table}[htp]
\centering
\caption{Công nghệ gắn vào sản phẩm.}
\label{tab:cong-nghe}
\begin{tabular}{p{3.4cm}p{3.8cm}p{6.6cm}}
\toprule
\textbf{Thứ dùng} & \textbf{Nhóm} & \textbf{Việc trong đồ án} \\
\midrule
React 19, TypeScript & Giao diện & SPA sau đăng nhập \\
Vite + nginx & Đóng gói FE & Dev proxy / production proxy \texttt{/api} \\
Tailwind CSS 4 & Trang trí & Không kéo bộ thành phần nặng \\
TanStack Query & Cache phía trình duyệt & Luật, hỏi trạng thái soát xét \\
Zustand & Phiên & Người dùng và khôi phục token \\
FastAPI & Máy chủ & REST + SSE bất đồng bộ \\
SQLAlchemy 2 + Alembic & ORM / migration & Schema có phiên bản \\
PostgreSQL 17 & CSDL & Nghiệp vụ + tìm chữ \\
ChromaDB & Vector & Cosine trên Điều \\
LangGraph & Điều phối & Bảy nút pipeline \\
slowapi & Giới hạn tần suất & Bảo vệ hạn mức mô hình \\
Docker Compose & Cụm & Postgres, backend, frontend \\
\bottomrule
\end{tabular}
\end{table}

\subsection{HTTP, JSON và CORS}

\boncau
{Cách trình duyệt nói chuyện với máy chủ: mỗi lần gọi độc lập, thân REST là
JSON, khi dev hai cổng khác nhau thì trình duyệt hỏi CORS~\cite{rfc9110}.}
{Toàn bộ nghiệp vụ đi qua HTTP. Body REST là JSON. Khi Vite (5173) gọi API
(8023) khác origin, trình duyệt bắt buộc có CORS.}
{GET đọc, POST tạo, PATCH sửa, DELETE xóa. Mã dùng: 200, 201, 202 (soát xét
nền), 400, 401, 403, 404 (kể cả xem chéo hội thoại), 422, 429, 503. CORS khóa
origin cấu hình, không \texttt{*}. Docker cùng origin nên CORS gần như không chạy.}
{Mọi request API. OPTIONS chỉ khi trình duyệt preflight. HTTPS/TLS không bật
trong Compose nội bộ --- khi đưa ra Internet phải đặt trước nginx.}

\subsection{Python}

\boncau
{Ngôn ngữ của toàn bộ tầng máy chủ: FastAPI, tách từ, LangGraph và SQLAlchemy
chung một codebase.}
{\texttt{underthesea} (tách từ BM25) ghim Python 3.10; FastAPI, SQLAlchemy,
LangGraph cùng một ngôn ngữ.}
{Quản lý gói bằng \texttt{uv}; type hint + Pydantic cho schema; \texttt{async}
cho I/O tới Postgres, Chroma và Gemini.}
{Toàn bộ tầng máy chủ. Không dùng cho frontend.}

\subsection{FastAPI}

\boncau
{Cửa HTTP duy nhất của SPA: nhận REST và đẩy SSE, kiểm body bằng Pydantic, tự
sinh mô tả OpenAPI~\cite{fastapi}.}
{Một request hỏi đáp chủ yếu chờ mạng: bất đồng bộ có lợi. SSE có
\texttt{StreamingResponse}. Schema khớp TypeScript phía trình duyệt.}
{Router tách theo nghiệp vụ (\texttt{auth}, hội thoại, hỏi đáp, luật, tài liệu,
lịch, quản trị). Pipeline RAG nằm tầng dịch vụ, không nằm trong router.}
{Mọi HTTP nghiệp vụ. Đường \texttt{/docs} dùng khi trình bày API, không phải
luồng người dùng DNNVV.}

\subsection{Pydantic}

\boncau
{Lớp biên giới API: body sai thì 422, không vào nghiệp vụ; schema khớp
TypeScript phía trình duyệt.}
{Lỗi 422 có cấu trúc khi client gửi sai trường; OpenAPI khớp schema thật.}
{Body đăng ký, đăng nhập, tin nhắn, kết quả soát xét, bộ câu chạy hàng loạt.}
{Mỗi lần FastAPI nhận hoặc trả JSON có khai báo mô hình.}

\subsection{SQLAlchemy 2.0 và Alembic}

\boncau
{Cách đọc ghi PostgreSQL bằng lớp Python, kèm nhật ký đổi cột để máy khác lên
cùng schema.}
{CRUD hội thoại, tài liệu, lịch không viết SQL lặp. Không sửa schema tay trên
máy rồi quên ghi.}
{Bản bất đồng bộ + \texttt{asyncpg}. User ứng dụng không \texttt{CREATE EXTENSION};
extension nằm \texttt{init.sql}.}
{Mọi đọc/ghi PostgreSQL nghiệp vụ. Vector không đi ORM --- nằm Chroma.}

\subsection{PostgreSQL}

\boncau
{Nơi giữ user, hội thoại, Điều chữ, lịch và ràng buộc (mã số thuế duy nhất, task
không trùng kỳ).}
{Cần ràng buộc tổ chức--user--hội thoại, FTS tiếng Việt tự dựng, và lịch tuân
thủ. Chroma không làm việc này.}
{Bản 17. Cổng host 23432. Corpus chữ và nghiệp vụ cùng một cụm; cột generated
\texttt{tsvector} sau \texttt{unaccent}.}
{Luôn chạy khi hệ thống sống. Tra cứu FTS không cần Gemini.}

\subsection{ChromaDB}

\boncau
{Nơi giữ vector từng Điều và tìm $k$ Điều gần câu hỏi theo cosine, chạy trong
cùng tiến trình backend.}
{Nhẹ, chạy kèm tiến trình backend, đủ \soDieu{} Điều; không thêm cụm Pinecone.}
{Volume \texttt{chroma\_db}. Manifest (mô hình, chiều, số bản ghi) so khớp lúc
khởi động. Đổi chiều thì xóa và nhúng lại.}
{Nạp/reindex và nút \texttt{retrieve}. Thiếu chỉ mục: hỏi đáp 503, FTS vẫn dùng.}

\subsection{LangGraph}

\boncau
{Cách nối bảy bước hỏi đáp thành một đồ thị có nhánh thoát sớm và có sự kiện
SSE từng nút.}
{Cần thoát sớm khi câu không phải pháp luật, và đẩy SSE theo từng nút. Không
dùng agent tự chọn công cụ tự do: thứ tự truy hồi rồi sinh là cố định.}
{Bảy nút; bộ nhớ ngắn nạp từ PostgreSQL, không dùng \texttt{InMemorySaver}.}
{Mỗi lượt hỏi đáp. Soát xét hợp đồng tái sử dụng hàm truy hồi, không đi đủ bảy
nút hội thoại.}

\subsection{Uvicorn}

\boncau
{Tiến trình chạy FastAPI trong Docker: một worker, cổng 8023, đủ để giữ SSE và
không đua file Chroma.}
{SSE cần vòng sự kiện bất đồng bộ; Gunicorn đồng bộ thuần không hợp.}
{Docker: một tiến trình, không reload, cổng 8023. Một worker để hai process
không đua file Chroma.}
{Luôn khi backend sống. Reload chỉ trên máy phát triển ngoài Compose.}

\subsection{Mô hình ngôn ngữ và Gemini}

\boncau
{Dịch vụ viết câu và nhúng chữ qua API (Gemini), không phải kho luật và không
được tinh chỉnh trọng số trong đồ án.}
{Không có GPU để tự host. Lớp tương thích OpenAI cho phép đổi \texttt{base\_url}
sang vLLM sau này mà không viết lại nút.}
{\texttt{gemini-2.5-flash}, nhiệt độ 0, tối đa 4096 token. Nhúng tách
\texttt{gemini-embedding-001}. Flash thay Pro: đủ câu ngắn trên $k=8$ Điều, rẻ
hơn trên HyDE + filter + generate.}
{Intent, HyDE, filter, generate, và nhúng. Thiếu khóa: health 200, chat 503.
Không fine-tune.}

\subsection{Underthesea}

\boncau
{Công cụ tách từ tiếng Việt cho BM25, để ``lao động'' không bị đếm thành hai
mảnh.}
{BM25 cần đơn vị từ. ``Lao động'' tách thành ``lao'' / ``động'' làm nhiễu IDF.}
{Nạp được thì dùng; không thì fallback regex. Cache theo dấu vân tay tokenizer.}
{Lúc dựng chỉ mục BM25 (khởi động hoặc cache lệch). Không dùng cho embedding.}

\subsection{React}

\boncau
{Thư viện vẽ toàn bộ màn hình sau đăng nhập; đổi trang không tải lại cả
site~\cite{react}.}
{Đề bài yêu cầu React. Ứng dụng sau đăng nhập, không cần SSR/SEO.}
{React 19. Tách \texttt{ChatPanel}, \texttt{CitationList}, \texttt{AgentTrace}.
Không dùng Next.js.}
{Mọi màn hình sau login. Trang login/register cũng là component React.}

\subsection{TypeScript}

\boncau
{JavaScript có kiểu: lệch tên trường API (token, citation) bị bắt lúc build.}
{Lệch \texttt{access\_token} / \texttt{accessToken} bắt lúc biên dịch.}
{Khai báo user, hội thoại, Điều, finding, sự kiện SSE. Khớp OpenAPI FastAPI.}
{Toàn bộ frontend.}

\subsection{Vite}

\boncau
{Công cụ dựng SPA: lúc lập trình có proxy \texttt{/api}, lúc đóng gói ra tệp
tĩnh cho nginx.}
{Khởi động nhanh; proxy \texttt{/api} giống nginx production.}
{\texttt{VITE\_API\_BASE\_URL} trống trong Docker: cùng origin.}
{Dev trên máy; image frontend chỉ chứa bản đã build.}

\subsection{Tailwind CSS}

\boncau
{Cách trang trí giao diện bằng class nhỏ trong JSX, không kéo bộ nút thương mại.}
{Không kéo bộ component thương mại; gói nhỏ, nền tối đủ đọc Điều dài.}
{Bản 4. Nút, ô nhập, thẻ tự viết.}
{Mọi trang giao diện.}

\subsection{TanStack Query}

\boncau
{Bộ nhớ tạm phía trình duyệt cho danh mục luật và trạng thái soát xét, tách
khỏi kho phiên đăng nhập.}
{Danh mục luật, cây Điều, poll soát xét --- không nhét vào store phiên.}
{Poll 2,5 giây khi soát xét. Tắt \texttt{refetchOnWindowFocus}.}
{Mọi GET danh sách/chi tiết. Không dùng cho luồng SSE (Fetch).}

\subsection{Zustand}

\boncau
{Kho nhỏ phía trình duyệt: chỉ user, token và hàm khôi phục phiên.}
{Chỉ cần user + token + \texttt{restore}; Redux thừa boilerplate.}
{Đọc token \texttt{localStorage} khi tải trang. Không chứa danh sách luật.}
{Khởi động SPA và mỗi lần login/logout.}

\subsection{React Router}

\boncau
{Cách đổi URL mà không tải lại trang; số hiệu có dấu \texttt{/} cần bắt phần
còn lại của đường dẫn.}
{Số hiệu văn bản chứa \texttt{/} cần splat \texttt{*}.}
{Route \texttt{/login}, \texttt{/chat/:id}, \texttt{/laws/:lawId/*},
\texttt{/contracts}, \texttt{/compliance}, \texttt{/admin/*},
\texttt{/admin/batch-eval}. Admin \texttt{lazy()}. \texttt{RequireAuth} chờ
restore xong.}
{Mỗi lần người dùng đổi trang.}

\subsection{Axios}

\boncau
{Client REST: gắn Bearer, làm mới token một lần khi 401, gỡ JSON khi gửi file.}
{Cần interceptor chung: Bearer, làm mới token một lần, hàng đợi 401.}
{Xóa \texttt{Content-Type} khi \texttt{FormData} để trình duyệt đặt multipart.}
{Mọi REST. Không dùng cho SSE --- hỏi đáp dùng Fetch.}

\subsection{nginx}

\boncau
{Cổng duy nhất người dùng mở: phát tệp Vite và chuyển \texttt{/api} vào
backend, tắt đệm cho SSE.}
{Trình duyệt chỉ mở một cổng; \texttt{/api} vào backend.}
{Image frontend phục vụ bản Vite. Host 5173.}
{Môi trường Docker / production. Dev dùng proxy Vite thay nginx.}

\subsection{Docker và Docker Compose}

\boncau
{Cách dựng cùng một cụm Postgres--API--nginx trên máy khác nhau, không lệch
phiên bản Python hay CSDL~\cite{docker}.}
{Cùng Python 3.10, Postgres 17, nginx trên máy khác nhau. Không lệch phiên bản
khi chấm đồ án.}
{Ba dịch vụ: \texttt{postgres} (healthcheck \texttt{pg\_isready}),
\texttt{backend} (đợi healthy, \texttt{/health}), \texttt{frontend}. Compose
chỉ đọc \texttt{.env} gốc, không đọc \texttt{backend/.env}.}
{Lúc dựng và vận hành cụm. Máy dev có thể chạy Vite + Postgres ngoài Compose.}

\subsection{JWT}

\boncau
{Hai loại token ký sẵn (access ngắn, refresh dài) để API biết ai đang gọi mà
không giữ phiên trên RAM~\cite{rfc7519}.}
{Không thêm Redis. Access ngắn hạn giảm cửa sổ lộ token.}
{HS256, \texttt{JWT\_SECRET\_KEY}. Claim \texttt{sub}, \texttt{type}
(\texttt{access}/\texttt{refresh}), \texttt{iat}, \texttt{exp}, \texttt{jti}.
Access 15 phút, refresh 7 ngày. Gọi API bằng refresh bị từ chối.}
{Mọi request sau login. Token ở \texttt{localStorage} --- còn rủi ro XSS; hướng
phát triển là cookie \texttt{HttpOnly}.}

\subsection{Argon2}

\boncau
{Cách lưu mật khẩu: băm argon2id, cố ý chậm, không lưu chữ gốc~\cite{argon2}.}
{SHA-256 quá nhanh, dễ dò GPU. Cần hàm cố ý chậm và tốn bộ nhớ.}
{Qua passlib. Không lưu plaintext. Hash hỏng không sập login --- verify trả sai.}
{Đăng ký và mỗi lần đăng nhập. Không dùng cho JWT.}

\subsection{Server-Sent Events}

\boncau
{Kênh máy chủ đẩy tiến độ và từng chữ xuống trình duyệt trong một kết nối HTTP,
không mở WebSocket.}
{Cần tiến độ từng nút và từng token. WebSocket thừa chiều client; polling chat
giật chữ.}
{Ba loại: \texttt{status}, \texttt{token}, \texttt{result}. Heartbeat tránh
proxy cắt. Client dùng Fetch (POST + Bearer) vì \texttt{EventSource} không gắn
header.}
{Mỗi lượt hỏi đáp. Soát xét hợp đồng dùng polling 2,5 giây, không SSE.}

\subsection{slowapi}

\boncau
{Chặn một tài khoản gọi chat hoặc tải file quá dày, để khỏi vét hạn mức mô hình.}
{Một tài khoản không được vét hạn mức Gemini.}
{Chat 20 lần/phút, tải tệp 30 lần/giờ, theo user (hoặc IP nếu chưa login).}
{Các endpoint tốn mô hình hoặc tốn đĩa. Tra cứu FTS không siết cùng mức chat.}

\subsection{Kiến trúc client--server và SPA}

\boncau
{Tách giao diện và máy chủ: trình duyệt tải một bộ JS, mọi nghiệp vụ đi
\texttt{/api/v1}.}
{Tách giao diện và máy chủ; trang sau login không cần HTML sẵn cho SEO.}
{Tiền tố \texttt{/api/v1}. Ba môi trường: Vite + Postgres local, Compose ba
container, tương lai tách domain.}
{Toàn bộ sản phẩm. Không dùng SSR.}

\section{Chi tiết từng nút pipeline LangGraph}

Mỗi nút đọc vài trường của \texttt{LegalAssistantState}, ghi vài trường, phát
SSE \texttt{status}. Prompt gốc không trích toàn văn (dài, dễ lệch khi sửa mã).

\subsection{Nút \texttt{analyze\_intent}}

\boncau
{Phân loại câu vào ô chat: cần tra luật hay trả lời ngắn.}
{Tránh nhúng và gọi retrieve cho ``xin chào'', ``cảm ơn'', câu ngoài miền.}
{LLM chỉ được in \texttt{NEXT} hoặc \texttt{SKIP}. Nhãn khác hoặc lỗi mạng
$\rightarrow$ ép \texttt{NEXT} (thà tra thừa). Bộ nhớ ngắn đã gắn từ PostgreSQL
trước khi gọi.}
{Mọi lượt chat. Trang chạy hàng loạt luôn \texttt{NEXT} để không câu nào bị SKIP.}

\subsection{Nút \texttt{prepare\_retrieval\_query}}

\boncau
{Tạo văn bản dùng để tìm: câu gốc, câu viết lại, hoặc đoạn HyDE.}
{Câu người dùng ngắn, giọng đời thường; văn luật khác từ vựng.}
{Sản phẩm: HyDE bật, rewrite tắt. Ưu tiên đoạn HyDE, rồi câu gốc. Lỗi một nhánh
thì dùng nhánh còn lại. Intent SKIP thì không gọi LLM.}
{Sau intent NEXT. Tắt/bật bằng \texttt{config.yaml}.}

\subsection{Nút \texttt{retrieve}}

\boncau
{Gọi \texttt{search\_legal\_articles}: Chroma + BM25, gộp RRF, lấy $k=8$.}
{Cần tập Điều cho filter và generate. Tìm đồng bộ chạy luồng phụ để heartbeat
SSE không bị nghẽn.}
{Lai, trọng số 3:1, $k_{\mathrm{RRF}}=60$. \texttt{search\_spaces} từ client chỉ
giao registry đã nạp.}
{Intent NEXT và registry sẵn. SKIP: danh sách rỗng. Ngoại lệ: retrieved rỗng,
generate nói thiếu căn cứ --- không bịa từ tham số mô hình.}

\subsection{Nút \texttt{rerank}}

\boncau
{Nút giữ chỗ cho xếp hạng chéo cặp (câu hỏi, Điều).}
{Cross-encoder chính xác hơn cosine nhưng đắt; Gemini không có cổng rerank.}
{Trên sản phẩm \textbf{tắt}: sao chép nguyên danh sách retrieve. Lỗi cũng
fallback, không nuốt ứng viên.}
{Chỉ khi tự host endpoint rerank và bật trong cấu hình.}

\subsection{Nút \texttt{llm\_filter}}

\boncau
{Hỏi PASS/DROP từng Điều: ``có giúp trả lời câu hỏi không?''}
{Thay rerank: giảm Điều lạc trước khi sinh câu.}
{Gọi song song, semaphore \texttt{max\_concurrency}. Item lỗi mạng được giữ
(PASS). Nếu PASS ít hơn \texttt{min\_keep} thì lấy thêm từ danh sách gốc.}
{Bật trên sản phẩm sau retrieve. Không chạy khi retrieved rỗng.}

\subsection{Nút \texttt{generate\_answer}}

\boncau
{Sinh câu tiếng Việt trên prompt hệ thống + Điều đã lọc + vài lượt gần nhất.}
{Đây là chỗ người dùng đọc kết luận; nhiệt độ 0 để diễn giải ổn định.}
{Tối đa 4096 token; stream SSE \texttt{token}. SKIP: không nhét Điều, không
citation. Tập rỗng: nhắc thừa nhận thiếu nguồn. Hậu kiểm citation nằm kề API,
không trong nút.}
{Sau filter (hoặc sau retrieve nếu filter tắt).}

\subsection{Nút \texttt{format\_submission}}

\boncau
{Gom \texttt{doc\_ref}, \texttt{article\_ref}, viện dẫn chéo đã lưu ở
\texttt{extra}.}
{UI cần danh sách căn cứ bấm được; lượt sau cần tin nhắn AI trong state.}
{Không gọi LLM. Không truy hồi thêm.}
{Luôn chạy cuối pipeline, kể cả nhánh SKIP (danh sách căn cứ rỗng).}

Bảng~\ref{tab:bay-nut} là cấu hình sản phẩm tại thời điểm báo cáo.

\begin{table}[htp]
\centering
\caption{Bảy nút LangGraph trên cấu hình sản phẩm.}
\label{tab:bay-nut}
\begin{tabular}{p{3.3cm}p{2.2cm}p{7.8cm}}
\toprule
\textbf{Nút} & \textbf{LLM?} & \textbf{Khi nào / ghi chú} \\
\midrule
\texttt{analyze\_intent} & Có & SKIP/NEXT; lỗi thì NEXT. \\
\texttt{prepare\_retrieval\_query} & Có (HyDE) & Rewrite tắt; HyDE bật. \\
\texttt{retrieve} & Nhúng câu hỏi & Hybrid RRF 3:1, $k=8$. \\
\texttt{rerank} & Không (tắt) & Giữ nguyên retrieved. \\
\texttt{llm\_filter} & Có & PASS/DROP; \texttt{min\_keep}. \\
\texttt{generate\_answer} & Có & Nhiệt độ 0; stream token. \\
\texttt{format\_submission} & Không & Chuẩn hoá ref, ghi memory. \\
\bottomrule
\end{tabular}
\end{table}

\section{So sánh phương án và lý do không chọn}

Các mục trước đã trả lời bốn câu cho cái đồ án \emph{dùng}. Mục này nêu cái đã
xem rồi bỏ: phương án là gì, vì sao không chọn, và khi nào sẽ xem lại.

\subsection{REST và GraphQL}

GraphQL cho client tự chọn trường, giảm overfetch. Đồ án schema ổn định, trang
nào cũng cần gần đủ trường Điều/citation; lợi GraphQL nhỏ, thêm endpoint
\texttt{/graphql} và cache phức tạp hơn. REST + OpenAPI khớp FastAPI có sẵn.

\subsection{SSE, WebSocket và polling}

WebSocket hai chiều, hợp chat typing đối phương. Đồ án chỉ server đẩy tiến độ
và token: SSE đủ, đi qua nginx đơn giản hơn WS. Polling cho soát xét (2,5 giây)
vì việc nền không cần stream token; polling chat thì lãng phí và giật chữ.

\subsection{JWT và phiên Redis}

Redis session dễ thu hồi. Đồ án không thêm dịch vụ Redis: JWT stateless, access
ngắn. Thu hồi sớm (khóa user) kiểm \texttt{is\_active} lúc mỗi request --- đủ
đồ án, chưa denylist \texttt{jti}.

\subsection{PostgreSQL và MongoDB}

Mongo linh hoạt JSON. Pháp lý cần ràng buộc, FTS, giao dịch đăng ký+tổ chức.
Postgres JSONB vẫn chứa \texttt{extra} và citation. Không cần hai CSDL.

\subsection{Chroma, FAISS và Pinecone}

FAISS mạnh khi tự host CPU/GPU, thiếu metadata store sẵn. Pinecone SaaS thêm
tài khoản và mạng. Chroma nhúng process, volume Docker, đủ \soDieu{} vector.

\subsection{Next.js và Vite SPA}

Next.js SSR/SEO. Mọi trang nghiệp vụ sau login: SEO không phải yêu cầu. Vite SPA
đơn giản, khớp đề bài React. Không dùng Next chỉ để ``thêm cho có''.

\subsection{Redux Toolkit và Zustand}

Redux hợp state lớn, middleware nhiều. Đồ án một store phiên: Zustand ít boilerplate.

\subsection{Celery/Redis và BackgroundTasks}

Hàng đợi thật chịu chết worker, retry, nhiều máy. Soát xét đồ án một container:
BackgroundTasks. Khi tệp dài và cần chắc chắn, chuyển Celery --- hướng Chương~5.

\subsection{Cookie HttpOnly và localStorage}

Cookie an toàn XSS hơn, thêm CSRF. localStorage nhanh cho SPA Axios. Đồ án chọn
cái sau, ghi nợ cái trước.

\subsection{Fine-tune và prompt}

Fine-tune tốn GPU, luật đổi phải train lại. Prompt + RAG cập nhật bằng nhập
JSON. Đồ án không fine-tune.


\section{Bảng tổng kết: vai trò -- thách thức -- hướng xử lý}

Ba bảng dưới đây chốt việc mỗi khối phải làm gì trong \emph{đồ án này}, chứ
không lặp lại định nghĩa ở trên.

\begin{table}[htp]
\centering
\caption{Tổng kết khối truy hồi lai.}
\label{tab:tong-ket-retrieval}
\begin{tabular}{p{2.2cm}p{11.5cm}}
\toprule
\textbf{Hạng mục} & \textbf{Nội dung} \\
\midrule
Vai trò & Đưa đúng tập Điều liên quan vào lời nhắc trước khi viết câu. \\
Thách thức & Lệch từ vựng người--luật; số hiệu dễ bị kênh vector làm tròn; hạn
mức nhúng. \\
Hướng xử lý & BM25 + Chroma + RRF 3:1; HyDE; nhúng lô 64 kèm nghỉ dần; cache
BM25. \\
\bottomrule
\end{tabular}
\end{table}

\begin{table}[htp]
\centering
\caption{Tổng kết khối sinh câu và hậu kiểm trích dẫn.}
\label{tab:tong-ket-generate}
\begin{tabular}{p{2.2cm}p{11.5cm}}
\toprule
\textbf{Hạng mục} & \textbf{Nội dung} \\
\midrule
Vai trò & Viết câu tiếng Việt có thể đối chiếu với Điều gốc. \\
Thách thức & Mô hình bịa số Điều; diễn giải sai trên Điều đúng; đốt hạn mức. \\
Hướng xử lý & Lời nhắc siết căn cứ; lọc PASS/DROP; gỡ citation ngoài tập tìm;
nhiệt độ 0; thiếu khóa thì 503, không sập cụm. \\
\bottomrule
\end{tabular}
\end{table}

\begin{table}[htp]
\centering
\caption{Tổng kết khối nghiệp vụ DNNVV (hợp đồng + lịch).}
\label{tab:tong-ket-sme}
\begin{tabular}{p{2.2cm}p{11.5cm}}
\toprule
\textbf{Hạng mục} & \textbf{Nội dung} \\
\midrule
Vai trò & Biến hỏi đáp thành việc kế toán và nhân sự làm hàng tuần. \\
Thách thức & PDF ảnh; hạn chót luật đổi; cách ly dữ liệu nhiều công ty. \\
Hướng xử lý & Cắt điều khoản rồi dùng lại hàm tìm; \soRule{} rule có
\texttt{legal\_refs}; cách ly theo tổ chức và người dùng; từ chối PDF không có
lớp chữ lúc tải. \\
\bottomrule
\end{tabular}
\end{table}

\section{Công trình liên quan}

Đồ án dựa trên RAG và truy hồi dày~\cite{lewis2020rag, karpukhin2020dpr}. BM25
và RRF là công cụ xếp hạng đã được công bố~\cite{robertson2009bm25, cormack2009rrf}.
HyDE được dùng để đưa câu hỏi gần giọng văn bản luật hơn trước khi
nhúng~\cite{gao2023hyde}. Công trình chatbot tuyển sinh~\cite{duong2026} cho
thấy RAG vận hành được trên dữ liệu hẹp tiếng Việt.

Điểm khác của đồ án: kho có cấu trúc Điều--Chương; tìm không dấu nhưng hiện còn
dấu; lịch tuân thủ gắn Điều trong kho; bốn nghiệp vụ trên một ứng dụng web đa
người dùng. Phần hiện thực ở Chương~3; phần chạy thử ở Chương~4.
